% !TEX root = <name of your master file>
%
% Chapter: Geometric Elegance versus Coordinate Practice in Classical
% Mechanics
%
% Written to be \input{} or \include{}d into a master document using
% the book (or report) class, alongside the companion chapter
% "Strange Things in Classical Mechanics" (label chap:strange-
% classical-mechanics), which this chapter cross-references once.
% If that chapter is not part of the same compilation the single
% \ref{chap:strange-classical-mechanics} below will simply come out
% as "??" -- harmless, and easily fixed by including both.
%
% Assumes the master preamble loads at least:
%   \usepackage{amsmath,amssymb,amsthm}
%   \usepackage{mathtools}
%   \usepackage[margin=1in]{geometry}   % optional

\chapter{Geometric Elegance versus Coordinate Practice in Classical
  Mechanics}
\label{chap:geometric-vs-coordinate}

\section{Introduction}

There are two ways to write down classical mechanics, and every
physicist learns them in the wrong order. First comes the practical
apparatus: pick generalized coordinates $q^i(t)$, write a Lagrangian
$L(q,\dot q,t)$ or a Hamiltonian $H(q,p,t)$, grind through the
Euler--Lagrange or Hamilton equations, solve. This is what one
actually computes with, and it is indispensable --- but it is also,
in a precise sense, an arbitrary choice of language: nothing about
the physics singles out these particular coordinates $q^i$ over some
other set. Only later, if at all, does a student meet the second
formulation: configuration space as an abstract smooth manifold,
velocities and momenta as elements of its tangent and cotangent
bundles, the dynamics generated by a connection or a symplectic
form defined without reference to any coordinate system whatsoever.

The two formulations describe exactly the same physical content; the
tension between them is not a tension between two theories but
between two languages for one theory. And as with any well-chosen
change of language, some things become obvious that were obscure,
and some things become laborious that were once immediate. This
chapter sets the two side by side, with a single worked example run
in both, and asks what each buys --- and costs.

\section{The Coordinate Form, and What It Hides}
\label{sec:coordinate-form}

The familiar machinery needs only a brief reminder. Generalized
coordinates $q^1,\dots,q^n$ parametrize the configurations of a
system; the Lagrangian $L(q,\dot q,t) = T - V$ is a function of
coordinates, velocities, and time; the equations of motion are
\begin{equation}
  \frac{d}{dt}\frac{\partial L}{\partial \dot q^i}
  - \frac{\partial L}{\partial q^i} \;=\; 0,
  \qquad i=1,\dots,n.
  \label{eq:el-coord}
\end{equation}
This is the form in which mechanics is actually solved: one grinds
out~\eqref{eq:el-coord} in whatever coordinates the problem hands
you, and out comes a trajectory.

Consider the simplest possible case: a free particle of mass $m$
moving in a plane, described in plane polar coordinates $(r,\phi)$
rather than Cartesian ones. The kinetic energy is $T =
\tfrac12 m(\dot r^2 + r^2\dot\phi^2)$, and there is no potential, so
$L=T$. Equation~\eqref{eq:el-coord} gives
\begin{equation}
  m\ddot r - m r\dot\phi^2 \;=\; 0,
  \qquad
  \frac{d}{dt}\bigl(m r^2\dot\phi\bigr) \;=\; 0.
  \label{eq:polar-free}
\end{equation}
The second equation is unremarkable --- it just says angular
momentum about the origin is conserved, as it must be for a free
particle. The first is the interesting one: it contains a term
$-mr\dot\phi^2$ that looks, for all the world, like an outward
centrifugal force. There is, of course, no force of any kind acting
on this particle; it travels in a perfectly straight line. The term
is not physics. It is the price of insisting on writing Newton's
second law in a coordinate system whose coordinate lines are
curved --- exactly analogous to the centrifugal and Coriolis terms
that appear when the same free-particle problem is worked in a
rotating frame, except that here the "frame" is not rotating at all;
only the \emph{coordinate grid} is curved.

This is not a quirk of polar coordinates specifically. In any
curvilinear coordinate system, extra terms of exactly this kind are
forced to appear, because the bare second derivative $\ddot q^i$ is
not, by itself, a coordinate-independent (tensorial) object; only the
combination
\begin{equation}
  \ddot q^i + \Gamma^i_{jk}(q)\,\dot q^j\dot q^k
  \label{eq:covariant-accel}
\end{equation}
transforms correctly under a change of coordinates, where
$\Gamma^i_{jk}$ are the Christoffel symbols of the coordinate system
(built from the kinetic-energy metric $ds^2 = \sum_i m_i\,dq^i
dq^i$, in the sense of Chapter~\ref{chap:strange-classical-mechanics}).
The Euler--Lagrange machinery automatically assembles the correctly
covariant combination~\eqref{eq:covariant-accel}, but it does so by
generating, coordinate by coordinate, extra terms that read exactly
like additional applied forces --- centrifugal, Coriolis, Euler ---
even when, as here, the true physics is completely force-free.

None of this is a defect of the coordinate approach; it is the
coordinate approach doing exactly what it is for. This is how
problems actually get integrated, in closed form or numerically, and
it is the only form in which mechanics ever makes contact with a
laboratory: no experiment outputs "a point of an abstract manifold,"
it outputs a number, which is to say a coordinate, read off a ruler,
a protractor, or a clock. A great deal of the technical skill of
mechanics --- finding action-angle variables for an integrable
system, Delaunay or Poincar\'e elements in celestial mechanics,
normal coordinates near an equilibrium --- consists precisely in
finding the coordinates in which a given problem's equations of
motion become as simple as possible. Coordinate choice is not
bookkeeping; it is frequently the whole content of solving a hard
problem.

But this leaves an uncomfortable question hanging. The same physical
system, worked in different coordinates, can produce equations of
motion that look wildly different in complexity, that display or
conceal conserved quantities, that contain "force" terms present in
one description and absent in another. Precisely because the
coordinate form does not, by itself, distinguish physics from
artifact, it cannot by itself answer a question that matters a great
deal: how much of what is written on the page in
equation~\eqref{eq:polar-free} is actually about the world, and how
much is only about the choice of $(r,\phi)$?

\section{The Geometric Reformulation}
\label{sec:geometric-form}

\subsection{Configuration space as a manifold}

The geometric answer to that question begins by refusing to
privilege any coordinate system in the first place. Configuration
space $Q$ is taken to be an abstract smooth manifold; a specific
choice of generalized coordinates $q^i$ is only ever a local chart on
$Q$, one among infinitely many equally valid alternatives, and
physical content is, by definition, whatever does not depend on
which chart one happens to be using.

Velocities of the system live in the tangent bundle $TQ$; the
Lagrangian is intrinsically a function $L : TQ\times\mathbb R \to
\mathbb R$, stated without reference to any coordinates at all. The
Euler--Lagrange equations~\eqref{eq:el-coord} themselves have a
coordinate-free meaning: they are the flow of a specific vector field
on $TQ$, picked out by the Poincar\'e--Cartan structure that $L$
induces on $TQ$, satisfying the second-order condition that this
flow really does generate velocities consistent with
positions~\cite{AbrahamMarsden1978}. What one writes down in
coordinates as~\eqref{eq:el-coord} is the local expression of this
single intrinsic object, valid in every chart, looking different in
each.

\subsection{Phase space and the symplectic form}

Passing (via the Legendre transform, where it is well defined) to
the momenta $p_i = \partial L/\partial \dot q^i$ carries the
description to the cotangent bundle $T^*Q$, the phase space. This
space carries a canonical structure that exists with reference to no
coordinates whatsoever: the \emph{tautological} (or Liouville) 1-form
$\theta$, defined at a point $(q,p)\in T^*Q$ by $\theta(v) =
p\bigl(\pi_*v\bigr)$ for any tangent vector $v$ there, where $\pi:
T^*Q\to Q$ is the projection. Its exterior derivative $\omega =
d\theta$ is the canonical symplectic form on $T^*Q$: closed,
nondegenerate, and built entirely out of the bundle structure of
$T^*Q$ over $Q$, with no reference to a coordinate system anywhere in
its definition.

Hamilton's equations are then the local expression of a single
intrinsic equation,
\begin{equation}
  \iota_{X_H}\,\omega \;=\; dH,
  \label{eq:hamiltonian-vf}
\end{equation}
which defines the Hamiltonian vector field $X_H$ for any smooth
function $H$ on phase space. In local coordinates $(q^i,p_i)$ in
which $\theta = p_i\,dq^i$ and $\omega$ takes its standard constant
form, \eqref{eq:hamiltonian-vf} reproduces exactly the familiar
\begin{equation}
  \dot q^i \;=\; \frac{\partial H}{\partial p_i},
  \qquad
  \dot p_i \;=\; -\,\frac{\partial H}{\partial q^i}.
  \label{eq:hamilton-coord}
\end{equation}
Coordinates in which $\omega$ takes this simple constant form are
called \emph{canonical}, or \emph{Darboux}, coordinates, after Gaston
Darboux, who proved in 1882 that such coordinates always exist, at
least locally, around every point of every symplectic
manifold~\cite{Darboux1882}. This is a striking fact with no analogue
in Riemannian geometry, where curvature is a genuine local invariant
obstructing the metric from being brought to a standard flat form
even at a single point's neighborhood. Symplectic manifolds have no
such local invariant: every symplectic manifold of a given dimension
looks, locally, exactly like every other one. This is precisely why
Hamilton's equations can always be brought to the same simple
form~\eqref{eq:hamilton-coord}, no matter how complicated the
physical system --- the complexity of a hard mechanics problem lives
entirely in the choice of \emph{which} canonical coordinates to use,
never in whether they exist.

\subsection{Canonical transformations reinterpreted}

In the coordinate treatment, canonical transformations tend to arrive
as a slightly mysterious piece of technology: a table of generating
functions $F_1,F_2,F_3,F_4$, and a set of conditions to check that a
proposed change of variables $(q,p)\mapsto(Q,P)$ preserves the form
of Hamilton's equations. From the geometric point of view, a
canonical transformation is simply a diffeomorphism of phase space
that preserves $\omega$ --- a \emph{symplectomorphism}. The
generating-function apparatus becomes a practical tool for
constructing particular symplectomorphisms, rather than the
definition of the concept. Darboux's theorem, in this light, is
exactly the statement that a symplectomorphism from a neighborhood
of any point to the flat model $(\mathbb R^{2n}, \sum dq^i\wedge
dp_i)$ always exists: "finding canonical coordinates" and "finding a
symplectomorphism to the standard model" are the same task, phrased
in two vocabularies.

\subsection{Geodesics, connections, and the return of the Christoffel symbols}

The same reinterpretation applies to the free-particle example of
Section~\ref{sec:coordinate-form}. Motion generated by a Riemannian
metric --- whether the plain kinetic-energy metric of a free system,
or the Jacobi metric that reformulates a potential as pure geometry,
as in Chapter~\ref{chap:strange-classical-mechanics} --- is
intrinsically \emph{geodesic} motion: $\nabla_{\dot x}\dot x = 0$,
where $\nabla$ is the Levi-Civita connection of the metric, an
object defined without reference to any coordinate chart. Written
out in coordinates, this intrinsic equation reads
\begin{equation}
  \ddot x^i + \Gamma^i_{jk}(x)\,\dot x^j \dot x^k \;=\; 0,
  \label{eq:geodesic-coord}
\end{equation}
with exactly the Christoffel symbols that appeared, unbidden, as
"fictitious centrifugal force" in equation~\eqref{eq:polar-free}.
That is not a coincidence or an analogy: a free particle moving in a
straight line, described in polar coordinates, \emph{is} geodesic
motion in the flat metric $ds^2 = dr^2 + r^2 d\phi^2$, and the
Christoffel symbols of that metric in polar coordinates are precisely
the coefficients that turned up as a spurious force term. "Fictitious
force" and "Christoffel symbol" are the same object seen from two
sides. One calls it a force only for as long as one has declined to
think geometrically about it.

\section{A Dictionary Between the Two Languages}

The following rough dictionary summarizes the correspondence between
the two ways of talking about the same mechanics.

\begin{center}
\begin{tabular}{@{}p{0.46\linewidth}p{0.46\linewidth}@{}}
\toprule
\textbf{Coordinate language} & \textbf{Geometric language} \\
\midrule
Generalized coordinates $q^i$ &
  Local chart on the manifold $Q$ \\[2pt]
Velocity $\dot q^i$ &
  Tangent vector $v\in T_qQ$ \\[2pt]
Canonical pair $(q^i,p_i)$ &
  Point of $T^*Q$, in Darboux coordinates \\[2pt]
Poisson bracket $\{f,g\}$ by partial derivatives &
  $\omega(X_f,X_g)$, coordinate-free \\[2pt]
Conserved quantity from a cyclic coordinate &
  Momentum map for a Lie-group symmetry of $\omega$ \\[2pt]
"Fictitious" centrifugal/Coriolis force &
  Christoffel symbol of the coordinate chart \\[2pt]
Coordinate singularity (poles of $(\theta,\phi)$, gimbal lock) &
  Ordinary, perfectly smooth point of the manifold \\[2pt]
Generating function $F_1,\dots,F_4$ &
  A specific symplectomorphism \\
\bottomrule
\end{tabular}
\end{center}

\section{What Each Form Buys You}

The geometric form earns its keep by telling you, in advance and for
free, which features of a set of equations of motion are physical
and which are artifacts of a chart. It explains why canonical
transformations exist and behave the way they do (Darboux's
theorem). It makes the relationship between symmetry and conservation
essentially automatic: a continuous symmetry is the flow of a vector
field preserving both $\omega$ and $H$, and when that symmetry comes
from the action of a Lie group, the associated conserved quantity is
the group's \emph{momentum map} --- the modern, coordinate-free
statement of Noether's theorem~\cite{Noether1918,MarsdenRatiu1999}.
It also exposes genuinely global, topological facts that no single
coordinate patch can see: the phase space of a rigid body, for
instance, is the cotangent bundle of the rotation group $SO(3)$,
which cannot be covered by a single nonsingular coordinate chart at
all --- the root cause of gimbal lock, which textbooks often present
as an engineering nuisance but which is, geometrically, an
unavoidable topological obstruction. A coordinate singularity, on
this view, is never a physical singularity: it is only ever the
failure of one particular chart, repaired by changing coordinates or,
better, by working with the coordinate-free structure directly, where
the question of a "singularity" never even arises.

The coordinate form earns its keep by being the only form in which a
problem is actually solved, in closed form or on a computer, and the
only form that connects directly to a measurement, which is always
reported as a number relative to some frame. Darboux's theorem
guarantees good coordinates exist near any point; it says nothing
whatsoever about how to find the ones adapted to a \emph{specific}
Hamiltonian's structure --- action-angle variables for an integrable
system, say, which is frequently the entire technical content of
actually solving a hard problem in mechanics. No amount of abstract
symplectic geometry does that work for you.

Seen from the perspective of the wider project of this book, the two
forms sit on opposite sides of its recurring divide: a coordinate
expression is what gets compared, number for number, against
measurement, and so belongs to the "real world" side of the ledger;
the manifold, the connection, the symplectic form are pieces of
"ideal world" mathematics, whose entire justification is that they
correctly organize and predict the coordinate-level facts, and
nothing more. Neither side is dispensable, and neither is more
fundamental than the other in any sense that would let one be
discarded.

\section{A Worked Contrast: The Spherical Pendulum}

A single example run in both languages makes the trade-off concrete.
A pendulum of length $\ell$ and bob mass $m$, free to swing in any
direction, has configuration space the sphere of radius $\ell$
centered on the pivot.

\emph{Coordinate form.} Using the polar angle $\theta$ from the
downward vertical and the azimuth $\phi$, the Lagrangian is
\begin{equation}
  L \;=\; \tfrac12 m\ell^2\bigl(\dot\theta^2 + \sin^2\theta\,
  \dot\phi^2\bigr) \;+\; mg\ell\cos\theta .
  \label{eq:sph-pendulum-L}
\end{equation}
The coordinate $\phi$ does not appear in $L$ itself, so its conjugate
momentum
\begin{equation}
  p_\phi \;=\; m\ell^2\sin^2\theta\,\dot\phi
  \label{eq:sph-pendulum-pphi}
\end{equation}
is conserved --- it is the angular momentum about the vertical axis
--- and the problem reduces to an effective one-dimensional motion in
$\theta$ alone, with effective potential $V_{\mathrm{eff}}(\theta) =
mg\ell\cos\theta + p_\phi^2/(2m\ell^2\sin^2\theta)$. This is the form
in which the problem is actually solved, e.g.\ by quadrature or
numerically.

\emph{Geometric form.} Configuration space is the sphere $S^2$
itself, with the circle group $SO(2)$ acting on it by rotation about
the vertical axis, an action that preserves the Lagrangian. The
conservation of $p_\phi$ is nothing but the momentum map for this
$SO(2)$ symmetry, an instance of Noether's theorem requiring no
mention of $\theta,\phi$ at all. And the coordinate singularities of
$(\theta,\phi)$ at the poles $\theta=0,\pi$ --- where $\phi$ becomes
undefined and equation~\eqref{eq:sph-pendulum-pphi} degenerates ---
are transparently artifacts of this particular chart: $S^2$ has no
distinguished points and no singularity there whatsoever. A pendulum
hanging straight down, or balanced straight up, is a perfectly smooth
physical configuration; it is only the coordinate system built from
$(\theta,\phi)$ that breaks down at exactly those two points, for the
same reason that all spherical coordinate systems break down at their
poles.

\section{Closing Remarks}

The two languages are not competitors; they specialize in different
questions. The geometric form is the place to ask what is really
happening, structurally, and what is guaranteed to remain true no
matter which coordinates one happens to choose. The coordinate form
is the place --- indeed the only place --- to ask what number a
system actually produces. A physicist who works only in coordinates
is at permanent risk of mistaking an artifact of a chosen chart for a
physical effect, exactly as happened, for a good part of the
twentieth century, with the coordinate singularity of the
Schwarzschild solution at its horizon: a singularity of the
particular $(t,r,\theta,\phi)$ coordinates originally used to write
the metric down, mistaken for decades for a genuine physical
singularity of spacetime itself, and only cleared up by passing to
coordinates (Eddington--Finkelstein, Kruskal) in which the horizon is
manifestly perfectly smooth --- a general-relativistic descendant, on
a much larger stage, of exactly the confusion embodied in the
centrifugal term of equation~\eqref{eq:polar-free}. A physicist who
works only in the geometric language, on the other hand, has an
elegant theory of everything and the means to compute nothing. The
working practice of mechanics has always required both, spoken
fluently, and switched between as the problem at hand demands.

\begin{thebibliography}{99}

\bibitem{AbrahamMarsden1978}
R.~Abraham and J.~E.~Marsden, \emph{Foundations of Mechanics}, 2nd
  ed.\ (Benjamin/Cummings, Reading, MA, 1978).

\bibitem{ArnoldMMCM1989}
V.~I.~Arnold, \emph{Mathematical Methods of Classical Mechanics}, 2nd
  ed., Graduate Texts in Mathematics 60 (Springer-Verlag, New York,
  1989).

\bibitem{MarsdenRatiu1999}
J.~E.~Marsden and T.~S.~Ratiu, \emph{Introduction to Mechanics and
  Symmetry}, 2nd ed., Texts in Applied Mathematics 17 (Springer, New
  York, 1999).

\bibitem{Darboux1882}
G.~Darboux, ``Sur le probl\`eme de Pfaff,'' \emph{Bull.\ Sci.\ Math.}
  (2) \textbf{6}, 14--36, 49--68 (1882).

\bibitem{Noether1918}
E.~Noether, ``Invariante Variationsprobleme,'' \emph{Nachr.\ Ges.\
  Wiss.\ G\"ottingen, Math.-phys.\ Kl.} 1918, 235--257.

\end{thebibliography}
